神经网络\label{term:1}是有史以来最美丽的编程范式之一. 在传统的编程方法中，我们告诉计算机要做什么，将大问题分解为许多计算机可以轻松执行的、精确定义的小任务. 相比之下，在神经网络中，我们不告诉计算机如何解决我们的问题. 相反，它从观测数据中学习，自行找出当前问题的解决方案.

从数据中自动学习听起来很有前景. 然而，直到 2006 年，除了少数专门问题之外，我们还不知道如何训练神经网络以超越更传统的方法. 2006 年发生的变化是发现了在所谓的深度神经网络中进行学习的技术. 这些技术现在被称为\emph{深度学习}（deep learning）\label{term:0}. 它们得到了进一步发展，如今深度神经网络和深度学习在计算机视觉、语音识别和自然语言处理中的许多重要问题上取得了杰出的性能. 它们正被 Google、Microsoft 和 Facebook 等公司大规模部署.

本书的目的是帮助你掌握神经网络的核心概念，包括深度学习的现代技术. 在读完本书后，你将编写出使用神经网络和深度学习来解决复杂模式识别问题的代码. 并且你将拥有一个基础，可以使用神经网络和深度学习来攻克你自己设计的问题.
\sectionTitle{基于原理的路线}{基于原理的路线}

本书的一个基本信念是，与其对一长串想法只有模糊的理解，不如扎实地掌握神经网络和深度学习的核心原理. 如果你很好地理解了核心思想，你就能迅速理解其他新材料. 用编程语言来类比，可以把它想象成掌握一门新语言的核心语法、库和数据结构. 你可能仍然只 ``知道'' 这门语言总量的一小部分——许多语言拥有庞大的标准库——但新的库和数据结构可以被快速而轻松地理解.

这意味着本书绝对不是一本如何使用某个特定神经网络库的教程. 如果你主要想学会使用某个库，那就不要读这本书！找到你想学的库，然后研读它的教程和文档. 但要注意. 虽然这样做能立即带来解决问题的回报，但如果你想理解神经网络中真正发生的事情，如果你想要的是多年后仍然有价值的洞见，那么仅仅学习某个热门库是不够的. 你需要理解神经网络工作方式背后那些持久、长存的洞见. 技术来来去去，但洞见是永恒的.
\sectionTitle{动手实践的路线}{动手实践的路线}

我们将通过攻克一个具体问题来学习神经网络和深度学习背后的核心原理：教计算机识别手写数字\label{term:2}的问题. 使用传统的编程方法，这个问题极其难以解决. 然而，正如我们将看到的，使用一个简单的神经网络，只需几十行代码，不需要任何特殊库，就能相当好地解决它. 更重要的是，我们将通过多次迭代来改进这个程序，逐步融入越来越多关于神经网络和深度学习的核心思想.

这种动手实践的路线意味着你需要一些编程经验才能阅读本书. 但你不需要成为一名专业程序员. 我用 Python（2.7 版）编写了代码，即使你不使用 Python 编程，只需稍加努力也应该很容易理解. 在本书的过程中，我们将开发一个小型神经网络库，你可以用它来实验和建立理解. 所有代码都可以在这里下载. 当你读完本书后，或者在阅读过程中，你可以轻松地选用一个功能更完整、面向生产环境的神经网络库.

与此相关的一点是，阅读本书对数学的要求并不高. 大多数章节中都有一些数学内容，但通常只是初等代数和函数图像，我预计大多数读者都能接受. 我偶尔会使用更高级的数学，但材料的组织方式使得即使某些数学细节让你感到困惑，你也能跟上. 唯一大量使用较深数学的章节是第 2 章，它需要一点多元微积分和线性代数. 如果你对这些不熟悉，我在第 2 章开头讨论了如何应对这些数学内容. 如果你觉得实在太吃力，可以直接跳到该章主要结果的总结. 无论如何，一开始都不需要为此担心.

一本书同时以原理为导向又注重动手实践，这是很少见的. 但我相信，如果我们把神经网络的基本思想构建出来，你会学得最好. 我们将开发活的代码，而不仅仅是抽象理论，是你可以探索和扩展的代码. 这样你就能从理论和实践两方面理解基本原理，并为进一步扩充你的知识做好充分准备.
